\section{Verified exercises, task sessions and rewards}
\subsection{What completion means}
Ascend distinguishes four kinds of evidence in challenge details. A bundled exercise
checks an exact result against a fresh, account-scoped server instance. A reasoning
case checks a scenario decision; it does not grade the accompanying explanation or
prove a hands-on assessment. A signed external observation checks a named game result
for an assigned instance. A file verifier checks an artifact's presence and size;
a person must still review the result, reasoning and transfer. None of these is a
professional certification. The supplied curricula remain unreviewed drafts.

In Tracks, choose \textbf{Browse learning tracks}, search by subject or role, inspect
entry requirements and install a track. Installation preserves the current mission.
Select a mission explicitly when ready. Today retains one next task and one optional
quick win; Path shows the level sequence, prerequisites and milestones.

Open a challenge and choose a realistic target before pressing \textbf{Start verified
exercise} or \textbf{Start timed task}. The default is the task's estimated minutes.
For a logging exercise, emit both success and failure, inspect the events for your
unique marker and enter the failed request ID and failure count. Other built-in
exercises cover IPv4 subnet calculations, positive and negative access decisions,
SHA-256 integrity, and login-event timelines. Scenario choices appear as readable
radio cards on phones. Wrong answers do not mark completion or award XP.

\subsection{Time and XP contract}
Time starts on the server when a task session starts. Merely opening a resource does
not start a timer. One session can run per account; complete or end it before
starting another. A task cannot be paused. Returning to its workspace preserves
the original target and running clock; changing an input does not reset either.
Closing the browser does not stop the server clock. A fresh exercise instance
changes its data without resetting the task session.

Let $t$ be elapsed seconds and $b$ the initial target seconds. Verified
completion awards
\[
 X(t,b)=\max\left(25,100-\left\lfloor50\frac{\max(0,t-b)}{b}\right\rfloor\right).
\]
Completion within target gives 100 XP; twice the target gives 50 XP; the floor is
25 XP. Only bundled exercises and signed CTF checks qualify. Files, self-marking
and skipped tasks receive no automatic skill XP. Start a session before an external
check to qualify. XP measures timed verification in this app, not independent speed
in an external game: the collector observes current game state, not the moment of
solving. The initial target is learner-controlled, so XP is personal feedback.

Each stable account/pack/task ID has at most one award. Completion or skipping ends
its session. Reopening, undoing a log or practicing again does not remove the historical
award or create another one. A finished session cannot be restarted as a new scored
attempt. Use fresh instances for unscored practice; never reset real history to farm XP.
\textbf{End session} stops timing without completing the task or awarding XP.
Verification after ending can record completion but cannot award this attempt XP.
Legacy paused rows remain readable for unscored practice and can be ended; current
clients cannot create or resume a paused task session. Mission pausing is separate:
it changes the selected learning track and does not pause a running task clock.

Starting opens a dedicated \texttt{\#challenge/PACK:NODE} workspace. The session strip
on every page has a \textbf{Resume challenge} action to return to the same exercise;
this action only navigates and does not restart timing. Reloading restores the server
instance and its event data. Draft answers are kept in this browser tab's session
storage and restored for that instance. They are cleared after successful submission
or an explicit fresh instance; they are not synchronized to another browser.

The page shows a synchronized countdown and reminders at 75 percent and the target.
The server checks every 30 seconds and can deliver the same two checkpoints through
the configured reminder channel when reminders are enabled. A suspended browser may
delay the display; notification delivery depends on the configured provider and
connectivity. These checkpoints are separate from the existing daily habit reminders.
No XP already earned is removed when a later session runs over time.

\subsection{CLI and API workflow}
All clients use the same application operations. Replace INSTANCE with the identifier
returned by start and use your actual generated answers:
\begin{verbatim}
ascend pack install verified-security-lab
ascend mission start verified-security-lab
ascend lab start verified-security-lab:lab-readiness --minutes 15
ascend lab show INSTANCE
ascend lab emit INSTANCE ok
ascend lab emit INSTANCE error
ascend lab submit INSTANCE '{"request_id":"YOUR_RESULT","errors":1}'
ascend session status
ascend session end
ascend session start PACK:NODE
ascend check PACK:NODE --complete
\end{verbatim}
Exercise start creates the timer and instance in one database transaction. Ordinary
tasks use \texttt{session start}. For an exercise already passed, submission
is idempotent. An unfinished instance expires after 24 hours; after ten failed
submissions, review the lesson and start a fresh instance. The server never evaluates
learner-supplied code. Stored expected results are not returned by the exercise API.
Native verification has a known timestamp and can create a dated evidence entry when
automatic evidence logging is enabled. External observations deliberately have no
solve timestamp; write a learning note to record today's practice and streak.
